\section{Degree-four linking correction in a two-truncated source}
\label{sec:136-two-truncated-linking-repair}

\paragraph{Archive and exact scope.}
This concerns cote136, physical pp.~117,119,121,
footers116,118,120, scans117,119,121.  The source in this
module is only a diagram on \(\Delta_{\leq2}^{\mathrm{op}}\).
All target theories are fully admissible theories restricted to
that diagram.  The absence of simplicial degree three is a
literal hypothesis, not an omitted coherence check.  The result
corrects degree-three graph linking through one full-boundary
degree-four Tate stage.  It does not construct an admissible
untruncated common source or an entire two-truncated iteration.

Let \(S=\Gamma_{\mathbb Z}(T)\), and put
\(J=\mathbb Q[T]/S\).  Use the injectivity of \(J\)
in nonnegative weight-graded \(S\)-modules proved in
ARCH-136-GRADED-INJECTIVES.  Let \(A\) be a two-truncated
strict differential PD algebra, with \(A_0=S\) and
\(A_1=F_1\), the specified free even zero-form unit-disk
algebra.  Suppose in every weight that
\[
 H^0(A_2)=S,\qquad H^1(A_2)=0,\qquad
              H^2(A_2),H^3(A_2)\text{ are integer torsion}.
                                                               \tag{136.TL1}
\]
Fix finitely many maps \(f_j:A\to(C_j)_{\leq2}\),
where the admissible \(C_j\) have the literal common first
skeleton \(F\) and these maps are the identity there.
Write \(G=\partial\Delta[2]=\operatorname{sk}_1\Delta[2]\),
oriented by the usual triangle boundary.

Assume that the degree-three graph-difference residues of the
old maps vanish in all weights \(w\geq W_0\), for every
finite modulus.  Equivalently, every old pair
\(dz=0,\ db=Nz\) in forms three and two in that range
has target primitive pairs with identical restrictions to
\(F(G)\).  Choose \(W_0\) also above the target vertical
and row bounds required for the graph comparisons below.
Those bounds involve only finitely many form degrees and are
independent of the modulus.

For any set of indices \(\alpha\), take old data in \(A_2\)
\[
 z_\alpha\in A_2^{v_\alpha,4},\quad
 e_\alpha\in A_2^{v_\alpha,3},\quad
 dz_\alpha=0,\quad de_\alpha=r_\alpha z_\alpha,
                         \qquad r_\alpha\geq1.
                                                               \tag{136.TL2}
\]
Prescribe complete old boundary data \(c_\alpha^\partial\)
in form three and \(t_\alpha^\partial\) in form two on
\(G\), satisfying the restrictions of the following equations.
Adjoin free cells along their full boundaries:
\[
 A\longrightarrow E,
 \qquad |c_\alpha|=3,\quad |t_\alpha|=2,
 \qquad dc_\alpha=z_\alpha,
 \qquad dt_\alpha=e_\alpha-r_\alpha c_\alpha.
                                                               \tag{136.TL3}
\]
Thus \(E_0=A_0\), \(E_1=A_1\), and \(E_2\) is
free as a graded PD algebra over \(A_2\) on these pairs;
its faces are the prescribed old sections.  Suppose strict
target extensions \(g_j:E\to(C_j)_{\leq2}\) have been
provided.  In the final paragraph of the proof we recall why
compatible old boundary data and extensions are available in
a sufficiently large fixed weight range.

Require every \(v_\alpha\) to exceed one fixed bound
\(V\), chosen so that in every target a coefficient
\(s\in S^{v_\alpha}\) is the relative period of a closed
two-form on \((\Delta[2],G)\), and its absolute second
cohomology vanishes.  The target relative comparison and row
bounds supply such a \(V\); no upper bound on the weights
or on the orders is assumed.

\paragraph{ARCH-136-TWO-TRUNCATED-LINKING-REPAIR.}
Under these hypotheses one can change only the new target
images \(g_j(t_\alpha)\) by closed two-forms with zero
full-boundary restriction, so that every degree-three
graph-difference residue of \(E\) vanishes in the same
uniform output tail, simultaneously for all finite moduli.
All old maps, all new \(c_\alpha\) images, and all
prescribed boundary data are retained.  The old second
cohomology is unchanged by the source attachment.

\begin{proof}
We give first the exact low-degree calculation.  There are no
new zeroforms or one-forms.  A possible new two-cycle has the
form \(a+\sum q_i t_i\), with \(q_i\in A_2^0\).
Its independent \(t_i\) and \(c_i\) coefficients imply
\(dq_i=0\) and \(r_iq_i=0\).  By (136.TL1),
\(q_i\in S\), so \(q_i=0\).  No one-form boundary
has changed.  Hence \(H^2(A_2)\to H^2(E_2)\) is an
isomorphism, and \(H^0,H^1\) are unchanged as well.

A three-cycle has precisely the form
\(a+\sum b_i t_i+\sum q_i c_i\), with \(b_i\) an
old one-form and \(q_i\) an old zeroform; products of two
new generators have form at least four.  Its equations are
\[
 db_i=0,\qquad dq_i+r_i b_i=0,\qquad
                da-\sum b_i e_i+\sum q_i z_i=0.
                                                               \tag{136.TL4}
\]
Write \(b_i=d\beta_i\) and subtract
\(d(\sum\beta_i t_i)\).  This gives a representative
\(a'+\sum s_i c_i\), with \(s_i=q_i+r_i\beta_i\in S\)
and \(da'+\sum s_i z_i=0\).  Boundaries change these
coefficients by \(-r_i\lambda_i\), \(\lambda_i\in S\).
An old three-cycle becomes a new boundary only if every such
coefficient is zero; torsion-freeness of \(S\) forces
\(\lambda_i=0\).  Therefore
\[
 \begin{gathered}
 H^2(E_2)=H^2(A_2),\qquad
 0\longrightarrow H^3(A_2)\longrightarrow H^3(E_2)
                              \longrightarrow K\longrightarrow0,\\
 L=\bigoplus_\alpha(S/r_\alpha)\langle v_\alpha\rangle,
 \qquad K=\ker\{L\longrightarrow H^4(A_2),\ c_\alpha\longmapsto[z_\alpha]\}.
 \end{gathered}                                             \tag{136.TL5}
\]
In particular the second and third cohomology of \(E_2\)
are integer torsion.  As two-truncated simplicial modules,
the relative coefficient modules \(L,K\) have normalized
components only in degree two.  All their boundary maps are zero.

We define the universal graph linking difference carefully.
For a torsion class \([x]\in H^3(E_2^w)\), choose
\(db=Nx\), with \(|b|=2\), and target primitives
\[
 da_j=g_jx,\qquad dq_j=g_jb-Na_j,
                  \qquad |a_j|=2,\quad |q_j|=1.
                                                               \tag{136.TL6}
\]
Fix target \(0\) as reference.  Since the maps agree on
the entire first skeleton, \((a_j-a_0)|_G\) is a closed
two-form in \(F(G)\).  Its second cohomology is zero
in the chosen uniform weight range; take
\(u_j\in F^{w,1}(G)\) with
\(du_j=(a_j-a_0)|_G\).  The one-form
\[
                   v_j=(q_j-q_0)|_G+Nu_j
                                                               \tag{136.TL7}
\]
is closed.  Let \(\operatorname{per}_G\) be the canonical
oriented isomorphism \(H^1(F^{w,*}(G))\simeq S^w\),
and define
\[
             \Delta_j([x])=
                  \frac{\operatorname{per}_G(v_j)}{N}\pmod S
                                     \quad\in J^w.
                                                               \tag{136.TL8}
\]
This is the degree-three counterpart of (136.LK10).

Changing \(u_j\) by a closed one-form changes (136.TL8)
by an integral period.  A change of \(q_j\) by a closed
target one-form is exact, so its graph period is zero.
Changing \(a_j\) by an exact two-form can be compensated
in \(q_j,u_j\), leaving (136.TL7) unchanged modulo exact
forms.  Target vanishing of \(H^1,H^2\) justifies these
statements and exhausts all target primitive ambiguities.

There is one additional issue compared with degree two:
changing \(b\) by a closed source two-form \(y\).
Choose target one-form primitives \(h_j\) of \(g_jy\).
Their difference on \(G\) is closed because all source
graph images are common.  The integral difference period
\[
                [y]\longmapsto
                    \operatorname{per}_G((h_j-h_0)|_G)
                            \quad\in S^w
                                                               \tag{136.TL9}
\]
is well-defined and additive: target primitive changes are
exact, and a source boundary may use its source primitive,
whose graph images agree.  It factors through \(H^2(E_2^w)\),
which is integer torsion by (136.TL5), whereas \(S^w\)
is torsion-free.  It is therefore zero.  This proves invariance
of (136.TL8) under the change of order relation.  It is precisely
the torsion-period mechanism of ARCH-136-TORSION-GRAPH-RESIDUE,
with the actual unchanged second cohomology identified here.

Changing the representative \(x\) by a source boundary is
handled by its source primitive in (136.TL6).  Passing to common
multiples of annihilators proves independence of \(N\),
additivity, and compatibility with all finite orders.  The
formulas also prove \(S\)-linearity.  Thus (136.TL8) is a
homomorphism on the indicated weight tail of \(H^3(E_2)\).
Its zero value is equivalent to vanishing of every finite-modulus
graph residue.  Explicitly, in the convention
\(d_{P_N}(a,b)=(da-Nb,-db)\), the difference of target
primitive pairs is \((q_j-q_0,-(a_j-a_0))\).
Subtracting the graph boundary \(d_{P_N}(0,u_j)\) reduces
it to \((v_j,0)\).  Its cone class is
\(\operatorname{per}_G(v_j)\pmod N\); the injection
\(S^w/NS^w\hookrightarrow J^w\) takes this to (136.TL8).

Old compatibility makes \(\Delta_j\) zero on \(H^3(A_2)\)
in the output tail.  By (136.TL5) it descends to a graded
map \(K_{\geq W_0}\to J\).  Graded injectivity extends
it to
\[
                            \lambda_j:L\longrightarrow J.
                                                               \tag{136.TL10}
\]
This is the entire normalized extension problem in the
two-truncated category: both source and target have only degree
two components.  There is no degree-three boundary subgroup
on which an additional chain condition could be required.
We extend from the whole output tail, so arbitrarily large
cell weights do not leave uncontrolled coefficient strips.

For each generator choose \(s_{j\alpha}\in S^{v_\alpha}\)
with
\(\lambda_j(c_\alpha)=s_{j\alpha}/r_\alpha\pmod S\).
This is possible because that value is killed by \(r_\alpha\).
The fixed bound \(V\) supplies a closed relative two-form
\(h_{j\alpha}\in C_j^{v_\alpha,2}(\Delta[2],G)\)
of relative period \(s_{j\alpha}\).  For the reference
target use zero.  Change
\[
                   g_j(t_\alpha)\longmapsto
                                 g_j(t_\alpha)+h_{j\alpha}.
                                                               \tag{136.TL11}
\]
Its differential and every boundary value are unchanged.
The actual free full-boundary pushout therefore gives strict
maps under \(A\), retaining all \(c_\alpha\) images.

To compute the sign, represent a three-class as
\(x=a+\sum_i s_i c_i\).  By the same common-multiple
argument as (136.OG12), enlarging \(N\) to kill the
remaining old torsion three-class if necessary, choose
\[
               db=Nx,\qquad
               b=b_{\mathrm{old}}-
                            \sum_i(N/r_i)s_i t_i.
                                                               \tag{136.TL12}
\]
The corrected image of \(x\) is unchanged.  Choose an
absolute one-form primitive \(H_{ji}\) of \(h_{ji}\)
in the target simplex.  Its graph restriction is closed and
has period \(s_{ji}\), by the relative connecting map.
The correction changes \(q_j\) in (136.TL6) by
\(-\sum_i(N/r_i)s_i H_{ji}\), while \(a_j,u_j\)
can be retained.  Consequently (136.TL8) changes by
\[
                    -\sum_i s_i\frac{s_{ji}}{r_i}
                       =-\lambda_j((s_i)_i).
                                                               \tag{136.TL13}
\]
This cancels the original linking map on \(K_{\geq W_0}\).
The independence and finite-modulus equivalence already proved
show that every desired graph residue vanishes, with the same
one bound and with no bound on the orders.

For completeness, the input boundary hypothesis has no hidden
degree-three simplex condition at this truncated stage.  In a
fixed sufficiently high weight, the graph algebra has
\(H^i(F(G))=0\) for \(i\geq2\).  Thus the restrictions
of the old pair (136.TL2) admit compatible \(c^\partial,t^\partial\)
in \(F(G)\).  Positive-form vertex restrictions are zero,
so these are genuine full-boundary data.  Choose arbitrary
target primitive pairs using vertical exactness through degree
four.  The difference from the prescribed graph data is a
closed degree-two Moore-cone cochain on \(G\), whose
cohomology is zero since \(H^2(F(G))=H^3(F(G))=0\).
Choose its graph cone primitive, extend its form-one and
form-two components to the target simplex by the fixed row
bounds, and subtract its differential.  This installs the
specified boundary data in every target.  All choices use one
fixed weight bound independent of the order and the number
of pairs.  The correction (136.TL11) then applies to these,
or to any other already prescribed such target extensions.
\end{proof}

\paragraph{Exact consequence and remaining gate.}
Full-boundary degree-four repairs can retain degree-three graph
linking in simplicial dimensions at most two, by a literal
coefficient-injective correction.  The untruncated counterexample
DEGREE-FOUR-FACTOR-OBSTRUCTION has a three-simplex and hence
does not satisfy this scope.  Completing a two-truncated
construction still requires preservation through its other
operations; passing from that truncated object to admissible
full theories requires a separate replacement theorem.

\paragraph{Verification status.}
The author completed a whole saved-proof reread, checking the
free low-degree calculation and its signs, unchanged second
cohomology, every choice in the universal graph linking map,
the exact torsion use in (136.TL9), all-modulus cone identification,
the nonnegative graded extension, the absence of a third
normalized component, the actual relative two-form correction,
and installation of complete boundary data at the stated fixed
bounds.  Root and reviewer cote123 subsequently completed whole
readings of the saved proof and reported PASS.  Their checks
include the low cohomology calculation without torsion-freeness
of the old zero-form group, the exact use of old second-cohomology
torsion, the all-modulus linking map, and the actual relative
two-form corrections with their signs.  No complete iteration
or full admissible source is asserted.
